\begin{frame}<1-2>[t]{新知探索}
\begin{campuscanvas}
% Source slide 9, objects 11; visible 1-2
\only<1-2|handout:1>{\node[anchor=north west,inner sep=0,text=black] at (70.000,150.000) {\begin{minipage}[t]{142.500mm}\raggedright\fontsize{16}{24.00}\selectfont \textbf{3. 元素的性质}\end{minipage}};}
% Source slide 9, objects 2,5; visible 1-2
\only<1-2|handout:1>{\node[anchor=north west,inner sep=0,text=black] at (70.000,245.000) {\begin{minipage}[t]{142.500mm}\raggedright\fontsize{12}{18.00}\selectfont \red{1. 确定性：}对于一个给定的集合，它的元素必须是确定的。也就是说，对于一个已知的集合来说，某个元素在不在这个集合里，是确定的，要么在，要么不在，不能含糊其辞。\end{minipage}};}
% Source slide 9, objects 4,6; visible 1-2
\only<1-2|handout:1>{\node[anchor=north west,inner sep=0,text=black] at (70.000,410.000) {\begin{minipage}[t]{142.500mm}\raggedright\fontsize{12}{18.00}\selectfont \red{2. 无序性：}集合中的元素无先后顺序。\end{minipage}};}
% Source slide 9, objects 7,8; visible 1-2
\only<1-2|handout:1>{\node[anchor=north west,inner sep=0,text=black] at (70.000,490.000) {\begin{minipage}[t]{142.500mm}\raggedright\fontsize{12}{18.00}\selectfont \red{3. 互异性：}集合中的元素必须互不相同。\end{minipage}};}
% Source slide 9, objects 9; visible 2
\only<2|handout:1>{\node[anchor=north west,inner sep=0,text=black] at (70.000,580.000) {\begin{minipage}[t]{142.500mm}\raggedright\fontsize{12}{18.00}\selectfont \textbf{集合相等：}只要构成两个集合的元素\red{一样}，我们就称这两个集合是\red{相等}的。\end{minipage}};}
\end{campuscanvas}
\end{frame}
